% !TeX program = pdflatex
% Una autodualidad centrada en el tritono: la clase de conjuntos 6-30 y su dual neo-riemanniano de orden 12.
% Versión en español (traducción con sentido del original en inglés, sixthirty_note.tex).
% Autor: Carles Marín (con Claude, Anthropic, como asistente de IA).
% Build: pdflatex sixthirty_note_es.tex (dos veces, por las referencias).

\documentclass[11pt]{article}
\usepackage[T1]{fontenc}
\usepackage[spanish,es-noindentfirst,es-noshorthands]{babel}
\usepackage{lmodern}
\usepackage{amsmath,amssymb,amsthm,booktabs,graphicx}
\usepackage[hidelinks]{hyperref}
\usepackage{xurl}
\usepackage[table]{xcolor}
\usepackage[margin=1in]{geometry}
\usepackage{microtype}
\usepackage{float}
\usepackage[section]{placeins}
\usepackage{tikz}
\usetikzlibrary{arrows.meta,positioning,calc}
\definecolor{ctA}{HTML}{1B6F8C}
\definecolor{ctB}{HTML}{B5530F}
\setlength{\emergencystretch}{2em}
\newtheorem{theorem}{Teorema}
\newtheorem{proposition}{Proposición}
\newtheorem{lemma}{Lema}
\newtheorem{remark}{Observación}
\newcommand{\lean}[1]{\texttt{#1}}
\newcommand{\Ztwelve}{\mathbb{Z}_{12}}
\newcommand{\Zn}[1]{\mathbb{Z}_{#1}}
\newcommand{\Dgrp}[1]{D_{#1}}
\newcommand{\Ahat}{\hat{A}}
\newcommand{\Stab}{\mathrm{Stab}}
\newcommand{\question}[1]{\smallskip\noindent\textit{#1}\smallskip}
% enlaces de audio por fichero (vivos en el repo público); en color de acento para que se vean clicables
\newcommand{\audiobase}{https://github.com/karlesmarin/music-math/raw/main/media}
\newcommand{\audio}[1]{\href{\audiobase/#1}{\textcolor{ctA}{\footnotesize\,[$\flat$\,audio]}}}

\title{\bf Una autodualidad centrada en el tritono:\\
  la clase de conjuntos 6-30 y su dual neo-riemanniano de orden 12}
\author{Carles Mar\'in\thanks{Preparado con la asistencia de un sistema de IA (Claude, Anthropic).
La matemática es clásica; toda afirmación, delimitación de alcance y limitación es responsabilidad del
autor. La frontera de la contribución se enuncia con claridad en la Sección~\ref{sec:novelty}, y es el
núcleo de Lean~4 ---no el asistente--- lo que certifica las partes formalizadas.}\\[3pt]
  \normalsize Investigador independiente\\
  \normalsize \texttt{karlesmarin@gmail.com}}
\date{Junio de 2026\\[3pt]\normalsize DOI: \href{https://doi.org/10.5281/zenodo.20820962}{10.5281/zenodo.20820962}
\quad\textbullet\quad \href{https://github.com/karlesmarin/music-math}{github.com/karlesmarin/music-math}}

\begin{document}
\maketitle

\begin{abstract}\noindent
El acorde de \emph{Petrushka} de Stravinsky superpone dos tríadas mayores separadas por un tritono.
Seguimos ese tritono hasta la teoría de grupos pura y de vuelta. Sobre el universo cromático
$\Ztwelve$, el grupo $T/I$ actúa de forma regular sobre la órbita de una \emph{clase de conjuntos} de
alturas ---admitiendo un grupo contextual dual simplemente transitivo de orden reducido, en el sentido
de Lewin--- precisamente cuando el estabilizador $T/I$ de la clase es \emph{normal}; entre las simetrías
de orden 2, esto singulariza al centro $\{T_0,T_6\}$, el tritono. La única clase no trivial en
$\Ztwelve$ con esa simetría es la \textbf{6-30} de Forte $=(013679)$, la clase de conjuntos del acorde
de Petrushka; su órbita tiene tamaño $12$ y su dual es diédrico de orden $12$ ($\cong\Dgrp{6}$). Damos la
caracterización con demostración, identificamos la instancia, enumeramos el fenómeno a través de los
universos (el conteo es $\mathrm{A}032239(n/2)$, brazaletes asimétricos de dos colores) y
\textbf{verificamos por máquina} en Lean~4 (sin \lean{sorry}, axioma-limpio) tanto el motor de dualidad
de fondo de Crans--Fiore--Satyendra como la estructura diédrica regular de orden 12 del dual de 6-30. La
matemática es clásica; las contribuciones son la caracterización específica, la instancia 6-30 con su
primer dual verificado por máquina, y la enumeración.
\end{abstract}

\section{El tritono, de Stravinsky a un grupo}\label{sec:intro}

En el segundo cuadro de \emph{Petrushka} (1911), Stravinsky hace sonar una tríada de Do mayor contra una
de Fa$\sharp$ mayor ---dos tríadas mayores separadas por un \emph{tritono}. Como clases de altura (con
$Do=0$) este ``acorde de Petrushka'' es $\{0,1,4,6,7,10\}$\,\audio{petrushka_chord.wav}, la clase de conjuntos
\textbf{6-30} de Forte,
un subconjunto de seis notas de la colección octatónica que la tradición rusa transmitió a Stravinsky a
través de Rimski-Kórsakov~\cite{Berger,Taruskin,vandenToorn}.\footnote{La lectura octatónica de
Stravinsky es la tradición analítica dominante (Berger, van den Toorn, Taruskin); para la posición
discrepante de que se aplica en exceso, véase Tymoczko~\cite{TymoczkoOct}. Usamos ``acorde de Petrushka''
como nombre de la sonoridad, por el ballet, sin atribuir el término a ningún autor concreto.}

\question{P\textsubscript{0}. ¿Por qué \emph{esta} colisión ---dos tríadas a un tritono--- suena tan
autocontenida? ¿Qué tiene de especial, estructuralmente, el tritono?}

Esta nota sigue esa pregunta hasta el álgebra de las transformaciones y de vuelta. La tradición
neo-riemanniana (Lewin~\cite{Lewin}; Crans--Fiore--Satyendra~\cite{CFS}) lee los acordes no por sus notas
sino por el grupo de operaciones que los relaciona, y empareja cada grupo de transformaciones con un
\emph{dual}. Preguntamos qué hace esa dualidad para un acorde \emph{simétrico} como el de Petrushka,
demostramos exactamente cuándo existe un dual reducido, y hallamos que la respuesta está gobernada en
todo momento por el tritono ---que resulta ser, a la vez, el corazón musical del acorde y el centro del
grupo. Toda afirmación sobre 6-30 que sigue está verificada por máquina en Lean~4.

\section{Preliminares: tríadas, el grupo $T/I$ y su dual}\label{sec:setup}

Para seguir el tritono necesitamos primero el reparto: los acordes, el grupo que los mueve y la noción de
grupo \emph{dual} de Lewin. Las clases de altura son $\Ztwelve$. El grupo $T/I$ es el grupo diédrico
$\Dgrp{12}$ de orden $24$: transposiciones $T_k:x\mapsto x+k$ e inversiones $I_k:x\mapsto -x+k$. Lo
identificamos con \lean{DihedralGroup 12} de Mathlib ($r\,i = T_i$, $sr\,j$ una inversión). Una
\emph{tríada} es un par $(\mathrm{fundamental},\mathrm{calidad})\in\Ztwelve\times\{\pm\}$ (mayor
$\{r,r+4,r+7\}$ o menor $\{r,r+3,r+7\}$); hay $24$, y $T/I$ actúa sobre ellas de forma simplemente
transitiva ---la representación regular por la izquierda de $\Dgrp{12}$. Las operaciones
neo-riemannianas $P,L,R$ generan el grupo dual de $T/I$ sobre las $24$ tríadas: dos subgrupos de
$\mathrm{Sym}(\Omega)$ son \emph{duales} (Lewin) cuando cada uno es el centralizador del otro y ambos
actúan de forma simplemente transitiva~\cite{CFS}.

\question{P\textsubscript{1}. La tríada es asimétrica, así que su órbita tiene $24$ formas y la dualidad
es de orden $24$. ¿Qué le pasa al dual cuando el objeto es un acorde \emph{simétrico}, cuya órbita es
menor?}

La dualidad es el caso regular del teorema clásico del centralizador (el centralizador de un grupo de
permutaciones regular es su representación regular opuesta; Dixon--Mortimer~\cite{DM} \S4.2, Teorema 4.2A;
Cameron~\cite{Cam}). Verificamos por máquina el motor y la instancia $T/I\!\leftrightarrow\!PLR$ en Lean
(Tabla~\ref{tab:lean}, espacio de nombres \lean{Duality}): el lema reutilizable del caso regular de
Wielandt \lean{centralizing\_fixedPoint\_eq\_one}, la cota \lean{card\_PLRgrp = 24}, y la dualidad misma
\lean{centralizer\_eq\_PLR}\,: $\mathrm{centralizador}(T/I)=\langle P,L,R\rangle$.

\section{¿Cuándo conserva un acorde simétrico su dual?}\label{sec:char}

Una clase simétrica tiene un estabilizador $T/I$ no trivial $H=\Stab(A)$, así que su órbita
$O=T/I\cdot A\cong\Dgrp{12}/H$ tiene tamaño $24/|H|<24$, y la acción $T/I$ sobre $O$ es transitiva pero ya
no libre. ¿Sigue $O$ portando un dual simplemente transitivo?

\begin{proposition}[Dualidad reducida]\label{prop:dual}
Sea $A$ un conjunto de clases de altura con estabilizador $T/I$ igual a $H=\Stab(A)\le\Dgrp{12}$ y órbita
$O=T/I\cdot A$. El centralizador de la imagen de $T/I$ en $\mathrm{Sym}(O)$ actúa de forma simplemente
transitiva ---un grupo dual de orden reducido en el sentido de Lewin--- \emph{si y solo si $H$ es normal
en $\Dgrp{12}$}, y entonces el dual tiene orden $|O|=24/|H|$.
\end{proposition}
\begin{proof}
$\Dgrp{12}$ actúa transitivamente sobre $O\cong\Dgrp{12}/H$ con estabilizadores de punto los conjugados
de $H$. Para una acción transitiva de un grupo $G$ con estabilizador de punto $H$, el centralizador $C$
de la imagen en $\mathrm{Sym}(O)$ es isomorfo a $N_G(H)/H$ y es semirregular; $C$ es transitivo ---
equivalentemente regular, equivalentemente el dual es simplemente transitivo--- si y solo si $N_G(H)=G$,
es decir, si y solo si $H\trianglelefteq G$ (\cite{DM} \S4.2; \cite{Cam}). Cuando $H\trianglelefteq G$,
$|C|=|G/H|=|O|$.
\end{proof}

\question{P\textsubscript{2}. ¿Cuál es esa simetría, musicalmente? ¿Existe un acorde cuyo estabilizador
sea exactamente el tritono ---y es único?}

Los subgrupos normales de $\Dgrp{12}$ contenidos en el grupo de rotaciones son exactamente los
$\langle T_d\rangle$; las reflexiones (inversiones) generan subgrupos no normales. Restringir a una
simetría \emph{de orden 2} aísla el tritono:

\begin{lemma}[El tritono es la involución central]\label{lem:tritone}
$T_6$ es el único elemento de orden $2$ del grupo de rotaciones $\Zn{12}$ (equivalentemente, $6$ es el
único punto fijo no nulo de $x\mapsto -x$), y $\{T_0,T_6\}=Z(\Dgrp{12})$ es el único subgrupo normal de
orden 2 de $\Dgrp{12}$. Una inversión aislada $\langle I_j\rangle$ es un subgrupo de orden 2 pero
\emph{no} es normal.
\end{lemma}

Así pues, entre las clases de conjuntos con una simetría de orden 2, la Proposición~\ref{prop:dual}
otorga un dual simplemente transitivo reducido (de orden 12) \emph{si y solo si} la simetría es el
tritono $T_6$. Una simetría de orden 2 no central (una inversión) no da dual. El barrido por fuerza bruta
de las $222$ clases (\lean{sage/duality\_sweep.sage}) corrobora el enunciado general sin excepciones; es
la demostración anterior la que lo convierte en teorema y no en mera observación.

\section{6-30: la respuesta, y es el acorde de Petrushka}\label{sec:sixthirty}

El Lema~\ref{lem:tritone} nos dijo dónde mirar: un acorde estabilizado por el tritono $T_6$ y nada más.
En $\Ztwelve$ hay exactamente uno, y ya lo conocemos.

\begin{theorem}[La instancia 6-30]\label{thm:sixthirty}
La única clase de conjuntos no trivial en $\Ztwelve$ cuyo estabilizador $T/I$ es exactamente
$\{T_0,T_6\}$ es la \textbf{6-30} $=(013679)$. Su órbita tiene tamaño $12$, y su grupo dual es regular y
diédrico de orden $12$ ($\cong\Dgrp{6}$).
\end{theorem}

$A=(013679)$ es la unión de los tres \emph{pares centrales} (tritonos) $\{0,6\},\{1,7\},\{3,9\}$, así que
es fijo por $T_6$ (una media vuelta del reloj de clases de altura,
Figura~\ref{fig:clock})\,\audio{t6_symmetry.wav}; al no tener simetría inversiva,
$\Stab(A)=\{T_0,T_6\}$ exactamente. Su inversión $(0,1,4,6,7,10)$ es el acorde de
Petrushka mismo (el par $6$-30A/B). \textbf{El ``ajá'' de P\textsubscript{0}:} la bitonalidad por tritono
que Stravinsky buscó de oído \emph{es} la simetría central $T_6$, y esa simetría central es justamente lo
que dota a 6-30 de su autodualidad de orden 12.

Todo este teorema está verificado por máquina en Lean~4 (archivo \lean{SixThirty.lean},
Tabla~\ref{tab:s30}), sin \lean{sorry} y axioma-limpio: $\Stab_T(A)=\{0,6\}$ y sin simetría inversiva
(\lean{stab\_T}, \lean{stab\_I\_empty}); órbita de tamaño $12$ (\lean{orbit\_card}); el núcleo de la
acción de $\Dgrp{12}$ sobre la órbita es el centro $\{T_0,T_6\}$ (\lean{Ker\_eq}), de modo que la imagen
es $\Dgrp{12}/Z$ de orden $12$ (\lean{card\_TIgrpO}); el dual (centralizador sobre la órbita) tiene orden
$12$ (\lean{card\_Cgrp}, construido como la \emph{representación regular opuesta}) y actúa
transitivamente, luego regularmente (\lean{Cgrp\_transitive}); y porta la huella diédrica
(\lean{dihedral\_fingerprint}: un generador $s$ de orden $6$, una involución $t\ne1$ con $tst=s^{-1}$ y
$st\ne ts$), fijándolo como diédrico de orden $12$.

\begin{figure}[h]
\centering
\begin{tikzpicture}[scale=1.9,every node/.style={font=\small}]
  \draw[thick] (0,0) circle (1);
  \foreach \a/\b in {0/6,1/7,3/9}{
    \draw[ctB,line width=1.3pt] ({90-30*\a}:1) -- ({90-30*\b}:1);
  }
  \foreach \k in {0,...,11}{
    \pgfmathtruncatemacro{\ina}{ ( \k==0 || \k==1 || \k==3 || \k==6 || \k==7 || \k==9 ) ? 1 : 0 }
    \ifnum\ina=1
      \filldraw[black] ({90-30*\k}:1) circle (0.05);
      \node at ({90-30*\k}:1.2) {\textbf{\k}};
    \else
      \draw[fill=white] ({90-30*\k}:1) circle (0.045);
      \node[gray] at ({90-30*\k}:1.2) {\k};
    \fi
  }
\end{tikzpicture}
\caption{La \textbf{6-30} $=\{0,1,3,6,7,9\}$ sobre el reloj de clases de altura (nodos rellenos). Es la
unión de los tres pares centrales $\{0,6\},\{1,7\},\{3,9\}$ (diámetros gruesos de tritono), luego fija
por $T_6$ (una rotación de $180^\circ$); con cero simetrías inversivas, \emph{no} es inversivamente
simétrica. Elegir tres de los seis pares de tritono, salvo rotación/reflexión, es la biyección con
brazaletes de la \S\ref{sec:enum} (Figura~\ref{fig:bracelet}).}
\label{fig:clock}
\end{figure}

\begin{figure}[h]
\centering
\begin{tikzpicture}[scale=1.0,every node/.style={font=\footnotesize},>={Stealth[length=2mm]}]
  \foreach \k in {0,...,5}{
    \coordinate (O\k) at ({90-60*\k}:2.3);
    \coordinate (I\k) at ({90-60*\k}:1.15);
  }
  \foreach \k in {0,...,5}{
    \pgfmathtruncatemacro{\next}{mod(\k+1,6)}
    \draw[ctA,->,thick] (O\k) -- (O\next);
    \draw[ctB,->,thick] (I\next) -- (I\k);
    \draw[dashed] (O\k) -- (I\k);
  }
  \foreach \k in {0,...,5}{
    \filldraw[fill=ctA!12,draw=ctA] (O\k) circle (0.30) node[black]{$T_{\k}A$};
    \filldraw[fill=white,draw=ctB] (I\k) circle (0.30) node[black]{$I_{\k}A$};
  }
\end{tikzpicture}
\caption{Grafo de Cayley del grupo dual $\Dgrp{6}$ de orden 12 actuando de forma simplemente transitiva
sobre la órbita $T/I$ de 12 elementos de 6-30: las seis formas por transposición $T_kA$ (exterior, con
$T_kA=T_{k+6}A$ identificadas) y las seis formas por inversión $I_kA$ (interior). Las flechas continuas
son un generador de orden 6 (la orientación se invierte en la clase lateral interior, como en todo grafo
de Cayley diédrico); las aristas discontinuas, un generador de orden 2. La regularidad de esta acción
(\lean{Cgrp\_transitive} + semirregularidad) es lo que hace del dual un grupo de orden $12$.}
\label{fig:cayley}
\end{figure}

\section{¿Con qué frecuencia ocurre esto?}\label{sec:enum}

\question{P\textsubscript{3}. ¿Es $\Ztwelve$ especial, o existe un autodual centrado en el tritono en
todo universo? ¿Cuántos hay?}

Para $n$ par, contamos las clases de conjuntos cuyo estabilizador $T/I$ es exactamente el centro
$\{T_0,T_{n/2}\}$ (dual $\cong\Dgrp{n/2}$, órbita $n$), recorriendo las uniones de los $n/2$ pares
centrales $\{x,x+n/2\}$ (\lean{sage/sixthirty.sage}):

\begin{table}[h]
\centering
\rowcolors{2}{ctA!7}{white}
\begin{tabular}{lrrrrrrrr}
\toprule
\rowcolor{ctA!22}
$n$    & 8 & 10 & \textbf{12} & 14 & 16 & 18 & 20 & 24\\
\midrule
conteo & 0 & 0  & \textbf{1}  & 2  & 6  & 14 & 30 & 127\\
\bottomrule
\end{tabular}
\caption{Clases autoduales centradas en el tritono por universo $n$: ninguna por debajo de $12$, única en
$12$ (la clase 6-30).}
\label{tab:counts}
\end{table}

Así, $\Ztwelve$ es el universo más pequeño donde el fenómeno existe, y allí es único. La sucesión
$1,2,6,14,30,62,127$ para $n=12,14,\dots,24$ es la \textbf{OEIS A032239} (``número de brazaletes
identidad (asimétricos) de $m$ cuentas, $2$ colores'') en $m=n/2$. La biyección
(Figura~\ref{fig:bracelet}): una clase autodual centrada en el tritono es una elección de subconjunto de
los $n/2$ pares centrales con grupo de simetría trivial, es decir, un brazalete asimétrico de dos colores
sobre $n/2$ cuentas. No hay fórmula cerrada simple (la tentadora $2^{m-5}-2$ ajusta $m=7..11$ y luego
falla en $m=12$: $\mathrm{A}032239(12)=127\ne126$); la enumeración es una especialización de
Burnside/Pólya. Lo nuestro es la caracterización musical, la biyección con brazaletes y la minimalidad de
6-30 en $\Ztwelve$.

\begin{figure}[h]
\centering
\begin{tikzpicture}[scale=1.4,every node/.style={font=\footnotesize}]
  \draw[gray!60] (0,0) circle (1);
  \foreach \i/\lab/\f in {0/{0,6}/1, 1/{1,7}/1, 2/{2,8}/0, 3/{3,9}/1, 4/{4,10}/0, 5/{5,11}/0}{
     \coordinate (b) at ({90-60*\i}:1);
     \ifnum\f=1 \filldraw[black] (b) circle (0.13);
     \else \draw[fill=white,thick] (b) circle (0.12); \fi
     \node at ({90-60*\i}:1.32) {$\{\lab\}$};
  }
\end{tikzpicture}
\caption{6-30 como brazalete de dos colores sobre $m=6$ cuentas: la cuenta $i$ es el par central
$\{i,i+6\}$, rellena si y solo si está en el conjunto. El patrón (relleno en $\{0,6\},\{1,7\},\{3,9\}$)
tiene simetría de rotación/reflexión trivial ---un brazalete \emph{asimétrico}, el objeto que cuenta
A032239; 6-30 es el único patrón de este tipo en $m=6$.}
\label{fig:bracelet}
\end{figure}

\section{Qué es el dual, musicalmente}\label{sec:music}

\question{P\textsubscript{4}. El dual es un grupo abstracto de orden $12$. ¿Para qué \emph{sirve}?}

Una acción de grupo simplemente transitiva es exactamente un \emph{Sistema de Intervalos Generalizado}
(GIS) en el sentido de Lewin~\cite[p.~157]{Lewin}: dados dos elementos cualesquiera, existe una operación
\emph{única} que lleva uno al otro. Así, el dual $\cong\Dgrp{6}$ es un GIS sobre las doce formas del
acorde de Petrushka\,\audio{orbit_12.wav} ---la Figura~\ref{fig:cayley} es su mapa. Cualquier forma del acorde alcanza
cualquier otra mediante un único movimiento dual; los doce colores de Petrushka se vuelven un único
espacio armónico navegable, el análogo para acordes simétricos de la red $P/L/R$ de las tríadas. La
Tabla~\ref{tab:family} sitúa el resultado en su familia.

\paragraph{Escúchalo.} Tres sonificaciones acompañan a esta nota ---el acorde de Petrushka (dos tríadas
a un tritono), la invariancia $T_6$ (6-30 y su transpuesta de tritono comparten las mismas seis clases
de altura) y un recorrido por las doce formas (la órbita del grupo dual, que pasa por el acorde de
Petrushka en $I_1$)--- como MIDI editable y WAV renderizado:
\url{https://github.com/karlesmarin/music-math/tree/main/media}.

\begin{table}[h]
\centering\small
\setlength{\tabcolsep}{5pt}
\rowcolors{2}{ctA!7}{white}
\begin{tabular}{@{}lcccc@{}}
\toprule
\rowcolor{ctA!22}
Objeto & Estabilizador & Órbita & Dual & Orden\\
\midrule
Tríada (Crans--Fiore--Satyendra) & $\{T_0\}$ & $24$ & $\Dgrp{12}$ & $24$\\
Ciclo hexatónico (Berry--Fiore) & $\{T_0,T_4,T_8\}$ (aum.) & $6$ & $S_3$ & $6$\\
\textbf{6-30 (esta nota)} & $\{T_0,T_6\}$ (tritono) & $12$ & $\Dgrp{6}$ & $12$\\
\bottomrule
\end{tabular}
\caption{Tres instancias de la misma maquinaria de dualidad por centralizador, según la simetría del
objeto.}
\label{tab:family}
\end{table}

\section{Qué está verificado por máquina}\label{sec:checked}

\question{P\textsubscript{5}. ¿Podemos fiarnos de todo esto?}

Todo lo afirmado sobre el motor de dualidad y sobre 6-30 está comprobado por el núcleo de Lean~4
(Tablas~\ref{tab:lean}--\ref{tab:s30}); \lean{\#print axioms} devuelve exactamente
\lean{[propext, Classical.choice, Quot.sound]} sin \lean{sorryAx} en cada teorema. La enumeración
(Tabla~\ref{tab:counts}) y la estructura nombrada $\Dgrp{6}$ se atestiguan además con scripts Sage/GAP
reejecutables (Tabla~\ref{tab:sage}).

\begin{table}[h]
\centering\small
\rowcolors{2}{ctA!7}{white}
\begin{tabular}{@{}l p{0.56\textwidth}@{}}
\toprule
\rowcolor{ctA!22}
\textbf{Teorema} & \textbf{Enunciado}\\
\midrule
\lean{centralizing\_fixedPoint\_eq\_one} & Caso regular de Wielandt (el centralizador de un grupo transitivo es semirregular).\\
\lean{PLR\_le\_centralizer}              & $\langle P,L,R\rangle\le\mathrm{centralizador}(T/I)$.\\
\lean{card\_PLRgrp}                      & $\mathrm{Nat.card}\,\mathrm{PLRgrp}=24$.\\
\lean{centralizer\_eq\_PLR}             & $\mathrm{centralizador}(T/I)=\langle P,L,R\rangle$ (la dualidad CFS).\\
\bottomrule
\end{tabular}
\caption{El motor de dualidad de fondo en \lean{NeoRiemannian.lean} (espacio de nombres \lean{Duality}),
verificado por máquina (Ladrillo 3).}
\label{tab:lean}
\end{table}

\begin{table}[h]
\centering\small
\rowcolors{2}{ctA!7}{white}
\begin{tabular}{@{}l p{0.60\textwidth}@{}}
\toprule
\rowcolor{ctA!22}
\textbf{Teorema} & \textbf{Enunciado}\\
\midrule
\lean{stab\_T}, \lean{stab\_I\_empty} & $\Stab_T(A)=\{0,6\}$; sin simetría inversiva.\\
\lean{orbit\_card}                    & la órbita $T/I$ de $A$ tiene $12$ elementos.\\
\lean{Ker\_eq}                        & el núcleo de la acción de $\Dgrp{12}$ sobre la órbita $=Z=\{T_0,T_6\}$.\\
\lean{card\_TIgrpO}                   & la imagen sobre la órbita $=\Dgrp{12}/Z$, orden $12$.\\
\lean{card\_Cgrp}                     & el dual (centralizador sobre la órbita) tiene orden $12$ (rep.\ regular opuesta).\\
\lean{Cgrp\_transitive}              & el dual es transitivo, luego regular (simplemente transitivo).\\
\lean{dihedral\_fingerprint}          & orden $12$; $s$ de orden $6$, $t^2=1$, $t\ne1$, $tst=s^{-1}$, $st\ne ts$.\\
\bottomrule
\end{tabular}
\caption{El dual diédrico regular de orden 12 de 6-30 en \lean{SixThirty.lean}, verificado por máquina
(Ladrillo 19) ---la primera demostración en Lean de la estructura del dual contextual de un acorde
simétrico.}
\label{tab:s30}
\end{table}

\begin{table}[h]
\centering\small
\rowcolors{2}{ctA!7}{white}
\begin{tabular}{@{}l p{0.58\textwidth}@{}}
\toprule
\rowcolor{ctA!22}
\textbf{Script} & \textbf{Qué atestigua}\\
\midrule
\lean{centralizer\_witness.sage} & GAP: $PLR==\mathrm{centralizador}$ exactamente; el dual de la 12-órbita de 6-30 $\cong\Dgrp{6}$.\\
\lean{sixthirty.sage}            & detalle de 6-30 (estab.\ $\{T_0,T_6\}$, órbita 12, dual orden 12) + conteos por $n$.\\
\lean{duality\_sweep.sage}       & las $222$ clases de $\Ztwelve$: regular $\iff$ autodual (0 excepciones).\\
\lean{sc630\_verify.sage}        & VI $=[2,2,4,2,2,3]$; estab.\ $\{T_0,T_6\}$; no inversivamente simétrica.\\
\bottomrule
\end{tabular}
\caption{Testigos Sage/GAP (Sage 10.8 + GAP), cada uno reejecutable.}
\label{tab:sage}
\end{table}

\section{Novedad, alcance y preguntas abiertas}\label{sec:novelty}

\paragraph{La contribución, dicha sin rodeos.} Esto es una nota de formalización $+$ caracterización, no
matemática nueva. El motor de dualidad es el caso regular del \textbf{teorema del centralizador de
Wielandt} (\cite{DM} \S4.2 Teor.\ 4.2A; \cite{Cam}); el marco de los grupos contextuales es
Fiore--Satyendra/Fiore--Noll~\cite{FS,FN11,FN25}. Lo nuestro: (i) la caracterización de la dualidad
reducida con demostración (Proposición~\ref{prop:dual}) y su corolario del tritono
(Lema~\ref{lem:tritone}); (ii) la instancia 6-30 con el \emph{primer} dual contextual de un acorde
simétrico verificado por máquina (Teorema~\ref{thm:sixthirty}, Ladrillo 19); (iii) la enumeración y la
biyección con brazaletes. Mathlib tiene \lean{DihedralGroup}, \lean{IsPretransitive} y el teorema de
Cayley, pero ningún resultado de centralizador-de-representación-regular; nuestro
\lean{centralizing\_fixedPoint\_eq\_one} y la construcción de la regular opuesta en \lean{SixThirty.lean}
cubren ese hueco como primeras formalizaciones.

\paragraph{Relación con Berry--Fiore.} Berry--Fiore~\cite{BF} demuestran que el grupo $PL$ de un ciclo
hexatónico es dual (en el sentido de Lewin) a su estabilizador $T/I$ ---la misma maquinaria. Los objetos
difieren (Tabla~\ref{tab:family}): el suyo es el estabilizador del eje aumentado $\{T_0,T_4,T_8\}$ con
dual $S_3$ de orden $6$; el nuestro es el centro de tritono $\{T_0,T_6\}$ con dual $\Dgrp{6}$ de orden
$12$. ``Tritono'', ``octatónica'' y ``6-30'' no aparecen en su artículo.

\paragraph{Relación con Peck.} Peck~\cite{Peck} generaliza los grupos conmutantes de Lewin a acciones
semirregulares y diagonales con ejemplos musicales; lo citamos por la maquinaria. Hasta donde sabemos su
tratamiento no incluye ni nuestra enumeración ni la instancia concreta de orden 12 ---verificado solo
contra el resumen y la literatura secundaria, pues el cuerpo del \emph{JMT} es de pago; los grupos
conmutantes octatónicos publicados en este linaje son de orden $8$ ($\cong\Dgrp{4}$, \cite{FN11}), no de
orden $12$. La instancia de orden 12 cae dentro del ámbito de los ``casos diagonales'' de Peck y se
matiza en consecuencia; la enumeración es la novedad mejor sostenida.

\paragraph{El linaje esquiva el caso (complementario, no ``olvidado'').} Lecturas de \cite{FS,FN11,FN25}
confirman que el autodual central de tritono no está enunciado: cada uno evita el régimen de órbita
reducida de forma distinta (pc-segmentos ordenados $+$ una condición de tritono; poda por hipótesis de
los recubrimientos; ``sin simetrías'' $+$ una condición de tritono que excluye $T_6$). Notablemente,
\cite{FN11} tabula $\{0,1,4,6,7,10\}$ (``Major Triad Tritone Mixture'') y estudia su recubrimiento por
$2$ tríadas (dual de orden $2$) ---la \emph{misma clase de conjuntos} que 6-30, sin formar nunca el dual
de orden 12 del hexacordo-como-punto. La literatura de asistentes de demostración solo tiene
Prismriver~\cite{Prism} (solo $T/I$, sin $PLR$/dualidad).

\question{P\textsubscript{6}. ¿Hacia dónde lleva el tritono después?}

\begin{itemize}\itemsep2pt
  \item \textbf{La familia general-$n$.} La Proposición~\ref{prop:dual} es independiente del universo;
        6-30 es la instancia no trivial mínima de una familia (dual $\cong\Dgrp{n/2}$ en el universo
        $n$). Un enunciado uniforme y su demostración en Lean quedan abiertos.
  \item \textbf{El isomorfismo nombrado.} Verificamos por máquina la \emph{huella} diédrica
        (Tabla~\ref{tab:s30}); el isomorfismo nombrado $\mathrm{dual}\cong\mathrm{DihedralGroup}\,6$
        requiere $\Dgrp{12}/Z\cong\Dgrp{6}$, hoy un hueco de Mathlib (aquí atestiguado por GAP).
  \item \textbf{Fourier.} La simetría $T_6$ es exactamente soporte en frecuencias pares de $\Ahat$ ---un
        hecho de magnitud; pero la autodualidad necesita información de \emph{fase} que $|\Ahat|$ no ve,
        la misma ceguera que la relación-Z. Una descripción de la autodualidad en el espacio de Fourier
        queda abierta.
  \item \textbf{Acordes mayores.} Análogos de séptima y de tetracordos (cf.\ \cite{FN25}): ¿hay un
        autodual centrado en el tritono entre ellos?
\end{itemize}

El tritono que abre el acorde de Petrushka resulta ser el centro de $\Dgrp{12}$, el obstáculo en el
espacio de Fourier y la semilla de una autodualidad de orden 12 ---un único objeto visto desde el
teclado, el reloj, el grupo y ahora el núcleo de Lean. Lo hemos seguido tan lejos como la máquina puede
certificar; señala más allá.

\section*{Disponibilidad de datos y código}
Las fuentes Lean (\lean{NeoRiemannian.lean}, \lean{SixThirty.lean}) y los scripts testigo Sage/GAP de las
Tablas~\ref{tab:lean}--\ref{tab:sage} acompañan a esta nota; cada teorema Lean es comprobable con
\lean{\#print axioms} y cada script es reejecutable.

\begin{thebibliography}{99}
\bibitem{CFS} A. S. Crans, T. M. Fiore, R. Satyendra, ``Musical actions of dihedral groups,''
  \emph{Amer. Math. Monthly} \textbf{116} (2009), no.~6, 479--495 (arXiv:0711.1873).
\bibitem{DM} J. D. Dixon, B. Mortimer, \emph{Permutation Groups}, Springer GTM 163, 1996
  (\S4.2, Teorema 4.2A).
\bibitem{Cam} P. J. Cameron, \emph{Permutation Groups}, London Math. Soc. Student Texts 45, CUP, 1999.
\bibitem{Lewin} D. Lewin, \emph{Generalized Musical Intervals and Transformations}, Yale Univ. Press,
  1987 (reimpr.\ Oxford Univ. Press, 2007); GIS $=$ acción simplemente transitiva, p.~157.
\bibitem{BF} C. Berry, T. M. Fiore, ``Hexatonic systems and dual groups in mathematical music theory,''
  \emph{Involve} \textbf{11} (2018), no.~2, 253--270 (arXiv:1602.02577).
\bibitem{Peck} R. W. Peck, ``Generalized commuting groups,'' \emph{J. Music Theory} \textbf{54} (2010),
  no.~2, 143--177. doi:10.1215/00222909-1214912.
\bibitem{FS} T. M. Fiore, R. Satyendra, ``Generalized contextual groups,'' \emph{Music Theory Online}
  \textbf{11}.3 (2005).
\bibitem{FN11} T. M. Fiore, T. Noll, ``Commuting groups and the topos of triads,'' in
  \emph{Mathematics and Computation in Music} (MCM 2011), LNCS \textbf{6726}, Springer, 2011, pp.~69--83
  (arXiv:1102.1496).
\bibitem{FN25} T. M. Fiore, T. Noll, E. Bonnell, H. Pyle, N. Rodriguez, M. Williams, S. Cannas,
  M. Andreatta, ``Transformations of triads and seventh chords: group extensions and duality,''
  arXiv:2507.20811 (2025).
\bibitem{Prism} L. Aniva, C. Wang, ``Prismriver: formalization of music theory and algorithmic
  composition in Lean~4,'' arXiv:2606.19936 (2026).
\bibitem{Berger} A. Berger, ``Problems of pitch organization in Stravinsky,'' \emph{Perspectives of New
  Music} \textbf{2} (1963), no.~1, 11--42.
\bibitem{Taruskin} R. Taruskin, \emph{Stravinsky and the Russian Traditions}, 2 vols., Univ. of
  California Press, 1996.
\bibitem{vandenToorn} P. C. van den Toorn, \emph{The Music of Igor Stravinsky}, Yale Univ. Press, 1983.
\bibitem{TymoczkoOct} D. Tymoczko, ``Stravinsky and the octatonic: a reconsideration,'' \emph{Music
  Theory Spectrum} \textbf{24} (2002), no.~1, 68--102.
\bibitem{Forte} A. Forte, \emph{The Structure of Atonal Music}, Yale Univ. Press, 1973.
\bibitem{OEIS} OEIS Foundation, sucesión A032239, ``Number of identity bracelets of $n$ beads of 2
  colors.''
\end{thebibliography}

\end{document}
